Large multimodal models (LMMs) have made impressive progresses and performed well on tasks such as image captioning, visual question answering, and grounded dialogue. Yet despite this progress, they continue to struggle with two fundamental challenges: learning new concepts beyond their training data, and reliably solving complex multimodal reasoning and generation tasks. Overcoming these limitations is essential if we want LMMs to function robustly in open-world settings, where new categories constantly emerge, and to support high-stakes applications like scientific analysis, education, or healthcare, which demand precise reasoning and generation. A crucial ingredient for overcoming these challenges is enabling models to reflect on and improve their own outputs. In this talk, I present three complementary efforts towards this vision. First, we explore how contrastive augmentation can expand models’ conceptual coverage, helping them recognize rare or fine-grained visual categories. Second, we introduce a multi-agent framework that decomposes complex multimodal generation into specialized roles, showing how structured collaboration improves reliability and scales with computational budget. Finally, we investigate fine-grained critique and correction in visual reasoning, proposing benchmarks and strategies that highlight reflection as a pathway to better learning and reasoning. Taken together, these directions sketch a roadmap towards self-improving multimodal models --systems that can adapt, reflect, and refine themselves in pursuit of deeper visual understanding and reasoning.
